\documentclass[11pt]{article}
\usepackage[T1]{fontenc}
\usepackage{lmodern}
\usepackage{amsmath,amssymb,amsthm,mathtools}
\usepackage[margin=1in]{geometry}
\usepackage{microtype}
\usepackage{xcolor}
\usepackage{graphicx}
\usepackage{hyperref}
\hypersetup{colorlinks=true,linkcolor=blue!55!black,citecolor=blue!55!black,urlcolor=blue!55!black}

\newtheorem{theorem}{Theorem}
\newtheorem{lemma}[theorem]{Lemma}
\newtheorem{remark}[theorem]{Remark}
\newcommand{\Om}{\Omega}
\newcommand{\pa}{\partial}
\newcommand{\LS}{\text{Lopatinskii--Shapiro}}
\newcommand{\Span}{\operatorname{span}}
\newcommand{\cD}{\mathcal D}

\title{The General-Boundary Spectral Inequality for the Bi-Laplacian}
\author{Candidate full solution packet (GPT5.6)}
\date{27 August 2026}

\begin{document}
\maketitle

\begin{center}
\fbox{\parbox{0.91\linewidth}{\textbf{Status.}
Candidate full solution, likely valid, subject to expert review. The proof
answers the spectral-inequality question in Section 1.6.2 of
arXiv:2109.01521v2. It does not address the neighboring questions on boundary
damping or quantitative boundary unique continuation.}}
\end{center}

\section{Source question and later theorem}

Let $\Om$ be a smooth bounded connected domain (or a compact smooth manifold
with boundary). Let $B_1,B_2$ be boundary differential operators of orders
$k_j$, with normal order at most three, and let
\[
  \mathcal P u=\Delta^2u,\qquad
  \cD(\mathcal P)=\{u\in H^4(\Om):B_1u=B_2u=0\text{ on }\pa\Om\}.
\]
The standing hypotheses in Le Rousseau--Zongo \cite{source} are:
\begin{enumerate}
\item $(\Delta^2,B_1,B_2)$ satisfies the $\LS$ condition;
\item $\mathcal P$ is self-adjoint and nonnegative on $L^2(\Om)$.
\end{enumerate}
Writing $(\phi_j)$ for an orthonormal eigenbasis and $(\mu_j)$ for the
eigenvalues, Section 1.6.2 asks whether, for every nonempty open
$\omega\subset\Om$, there is $C>0$ such that
\begin{equation}\label{goal}
 \|u\|_{L^2(\Om)}\leq C e^{C\mu^{1/4}}\|u\|_{L^2(\omega)},
 \qquad u\in\Span\{\phi_j:\mu_j\leq\mu\}.
\end{equation}
The exact source passage is on page 8:
\begin{center}
\includegraphics[width=0.96\linewidth]{figures/open_problem_crop.png}
\end{center}

Zongo--Robbiano \cite[Theorem 1.3, p.~6]{support} prove \eqref{goal} under
the preceding two hypotheses and the additional $\LS$ condition for
\[
 Q=D_s^4+\Delta^2
\]
with the same boundary operators on $(0,S_0)\times\pa\Om$. The source paper
had isolated precisely this extra condition as the obstruction: ordinary
$\LS$ need not imply parameter $\LS$ without further structure
\cite[p.~8]{source}. The point below is that the already assumed
self-adjoint nonnegativity rules out the obstruction.

\section{Proof intuition}

The augmented boundary condition is a negative-resolvent condition in
disguise. After Fourier transform in the added variable $s$, failure at a
frequency $\sigma\ne0$ produces a nonzero exponentially decaying solution of
the frozen normal problem
\[
   (P_0(\xi',D_t)+\sigma^4)w=0,\qquad B_{j,0}(\xi',D_t)w(0)=0.
\]
Thus the frozen half-line realization has the negative eigenvalue
$-\sigma^4$. Localizing this mode at spatial scale $h$ gives a boundary
packet with Rayleigh quotient $-\sigma^4h^{-4}+o(h^{-4})$. The principal
boundary traces vanish; the original elliptic boundary condition supplies a
trace correction smaller by $h^{1/2}$. The corrected packets lie in
$\cD(\mathcal P)$, contradicting $\mathcal P\ge0$.

This proves the augmented $\LS$ condition automatically. The spectral
inequality is then exactly Zongo--Robbiano's theorem.

\section{The automatic parameter condition}

We record the standard boundary-layer implication in the precise form needed
here. Principal symbols are taken in boundary normal coordinates at a fixed
$x^0\in\pa\Om$; $t\ge0$ is the inward normal variable.

\begin{lemma}[Trace correction at boundary-layer scale]\label{trace}
Assume $(P,B_1,B_2)$ is a normal fourth-order elliptic boundary problem of the
source class, satisfying ordinary $\LS$, whose realization is self-adjoint.
Suppose $w(t)$ is exponentially decreasing and
its principal boundary traces obey
\[
 B_{j,0}(x^0,\xi',D_t)w(0)=0,\qquad j=1,2.
\]
Let $u_h$ be the normalized packet
\begin{equation}\label{packet}
 u_h(x',x_d)=h^{-(d+1)/4}\chi(x'/\sqrt h)
 e^{ix'\cdot\xi'/h}\chi_0(x_d)w(x_d/h)
\end{equation}
in a boundary chart, where the cutoffs equal one near the origin. Then there
are smooth $v_h$, supported in a fixed slightly larger chart, such that
\[
 B_j(u_h+v_h)=0,\qquad
 \|v_h\|_{L^2}=O(h^{1/2}),\qquad
 \|Pv_h\|_{L^2}=O(h^{-7/2}).
\]
\end{lemma}

\begin{proof}
The normalization in \eqref{packet} gives $\|u_h\|_{L^2}=\|\chi\|_2
\|w\|_2+o(1)$ after a harmless fixed rescaling. A derivative on the carrier
or on $w(x_d/h)$ costs $h^{-1}$, whereas a derivative on the tangential
cutoff costs only $h^{-1/2}$. Freezing coefficients at $x^0$ therefore leaves
a relative $O(h^{1/2})$ error. Since the principal $B_j$ traces vanish, this
gives
\begin{equation}\label{trace-size}
 \|B_ju_h\|_{H^{7/2-k_j}(\pa\Om)}=O(h^{-7/2}),
\end{equation}
one half-power below the $H^4$ size $O(h^{-4})$ of the packet.

For completeness, use the boundary jet map
$J=(\gamma_0,\gamma_1,\gamma_2,\gamma_3)$ and write each $B_j$ as a
tangential differential matrix acting on $J$. On the $h^{-1}$ frequency
annulus, the normal elliptic boundary system gives a uniformly bounded right
inverse for the rescaled two-row boundary-symbol matrix. For $\xi'\ne0$ this
is the ordinary $\LS$ right inverse. At the pure normal ray $\xi'=0$, it is
the normality of the two independent boundary conditions; loss of rank there
would leave a nonzero unconstrained Green boundary jet and is incompatible
with a self-adjoint fourth-order realization. Apply the right inverse to
$-h^{k_j}B_ju_h$, and lift the resulting four jets by the usual Poisson trace
extension $e^{-x_d/h}$ (with cutoffs). The rescaled estimates from
\eqref{trace-size} are
\[
 \|v_h\|_{L^2}=O(h^{1/2}),\qquad
 \|D^\alpha v_h\|_{L^2}=O(h^{1/2-|\alpha|}),\quad |\alpha|\le4.
\]
Lower-order coefficient and cutoff errors are removed successively; the
resulting geometric symbol expansion has the same bounds. Equivalently, this
is the standard semiclassical right inverse for a normal elliptic boundary
trace system, as in the parameter boundary calculus of
Agranovich--Vishik \cite{AV}. A final smoothing correction makes the boundary
conditions exact without changing the estimates. Thus
$B_j(u_h+v_h)=0$, and the last display gives the claimed estimate for $Pv_h$.
\end{proof}

\begin{lemma}[Nonnegativity forces augmented $\LS$]\label{automatic}
Under the two standing source hypotheses, $(D_s^4+\Delta^2,B_1,B_2)$
satisfies the $\LS$ condition at every nonzero tangential covector
$(\sigma,\xi')$ on $(0,S_0)\times\pa\Om$.
\end{lemma}

\begin{proof}
At $\sigma=0$ and $\xi'\ne0$, this is the ordinary source $\LS$ condition.
Suppose instead that it fails at $(x^0,\sigma,\xi')$ with $\sigma\ne0$.
By the normal-root formulation of the complementing condition, there is a
nonzero exponentially decreasing $w$ on $\mathbb R_+$ satisfying
\begin{equation}\label{normalmode}
 (P_0(x^0,\xi',D_t)+\sigma^4)w=0,\qquad
 B_{j,0}(x^0,\xi',D_t)w(0)=0.
\end{equation}
Normalize the packet \eqref{packet} so that $\|u_h\|_2=1+o(1)$. The frozen
equation, coefficient variation on the $\sqrt h$ tangential scale, and cutoff
counting give
\begin{equation}\label{quasi1}
 \|(P+h^{-4}\sigma^4)u_h\|_2=O(h^{-7/2}).
\end{equation}
Let $v_h$ be supplied by Lemma \ref{trace}, and put $q_h=u_h+v_h$. Then
$q_h\in\cD(\mathcal P)$, $\|q_h\|_2=1+o(1)$, and
\begin{equation}\label{quasi2}
 \|(\mathcal P+h^{-4}\sigma^4)q_h\|_2=O(h^{-7/2}),
\end{equation}
because $h^{-4}\|v_h\|_2=O(h^{-7/2})$.

Self-adjoint nonnegativity and \eqref{quasi2} now imply
\begin{align*}
 0&\le \langle\mathcal Pq_h,q_h\rangle\\
 &=-h^{-4}\sigma^4\|q_h\|_2^2
   +O(h^{-7/2})\|q_h\|_2,
\end{align*}
which is negative for small $h$. This contradiction proves the augmented
$\LS$ condition.
\end{proof}

\section{Full answer}

\begin{theorem}[General-boundary bi-Laplacian spectral inequality]\label{main}
Assume the boundary operators and the realization $\mathcal P$ satisfy the
standing hypotheses of Le Rousseau--Zongo \cite{source}: ordinary $\LS$,
self-adjointness, and nonnegativity. For every nonempty open
$\omega\subset\Om$ there is $C>0$ such that, for all $\mu\ge0$,
\[
 \|u\|_{L^2(\Om)}\le C e^{C\mu^{1/4}}\|u\|_{L^2(\omega)},
 \qquad u\in\Span\{\phi_j:\mu_j\le\mu\}.
\]
Thus the spectral-inequality question in Section 1.6.2 of
arXiv:2109.01521v2 has an affirmative answer.
\end{theorem}

\begin{proof}
Lemma \ref{automatic} supplies the only additional hypothesis in
Zongo--Robbiano \cite[Theorem 1.3]{support}. Applying that theorem gives the
displayed inequality with the exponent $\mu^{1/4}$.
\end{proof}

\section{Verification, scope, and novelty check}

\paragraph{Hypothesis match.}
The source question occurs on page 8 of arXiv:2109.01521v2. The supporting
spectral theorem occurs as Theorem 1.3 on page 6 of arXiv:2309.02181v1. Its
extra assumption is exactly ordinary $\LS$ for the augmented operator
$D_s^4+\Delta^2$. Lemma \ref{automatic} derives that assumption from the
source's self-adjoint nonnegativity.

\paragraph{Sanity checks.}
The sign is forced: Fourier transformation in $s$ changes $D_s^4$ into
$+\sigma^4$, so a failed augmented condition gives a frozen mode at the
negative energy $-\sigma^4$. The packet error is $h^{-7/2}$ versus the main
$h^{-4}$ term, leaving the required half-power margin. When $\sigma=0$, the
argument correctly reduces to the source $\LS$ condition.

\paragraph{Bounded novelty search.}
On 27 August 2026 we searched the run indexes, arXiv/web results for the exact
phrases ``spectral inequality'', ``general boundary conditions'',
``bi-Laplace'', and ``Lopatinski--Shapiro'', and the citation trail of
arXiv:2109.01521. We found arXiv:2309.02181, which proves the desired estimate
with the additional augmented $\LS$ hypothesis and cites the source paper,
but no primary source removing that hypothesis by self-adjoint nonnegativity.
This is a bounded search, not a claim of exhaustive novelty.

\paragraph{Human review recommendation.}
The result is recommended for expert review as a likely valid full solution.
The decisive review point is the standard high-frequency right inverse in
Lemma \ref{trace}; a reviewer should check its normality convention against
the boundary-operator class in the source. No claim is made about Section
1.6.1 (boundary damping) or Section 1.6.3 (a quantitative interpolation
inequality).

\begin{thebibliography}{9}
\bibitem{source}
J. Le Rousseau and E. Wend-Benedo Zongo,
\emph{Stabilization of the damped plate equation under general boundary
conditions}, arXiv:2109.01521v2 (2021); J. \`Ec. polytech. Math. 10 (2023),
1--65.

\bibitem{support}
E. Wend-Benedo Zongo and L. Robbiano,
\emph{Null-controllability for a fourth order parabolic equation under
general boundary conditions}, arXiv:2309.02181v1 (2023), especially Theorem
1.3.

\bibitem{AV}
M. S. Agranovich and M. I. Vishik,
\emph{Elliptic problems with a parameter and parabolic problems of general
type}, Russian Math. Surveys 19:3 (1964), 53--157.
\end{thebibliography}

\end{document}
